% 3.6.2 定性示意图；不基于任何具体实验温度采样。
\begin{figure}[htbp]
\centering
\resizebox{0.91\textwidth}{!}{%
\begin{tikzpicture}[font=\small,>=Latex]
 \fill[red!7] (.90,2.6) rectangle (10.3,3.90);
 \draw[->,thick] (.90,.45)--(10.65,.45) node[right] {时间};
 \draw[->,thick] (.90,.45)--(.90,4.15) node[above] {温度};
 \draw[densely dashed,red!65!black,thick] (.90,2.6)--(10.25,2.6);
 \node[font=\scriptsize,red!65!black,anchor=west] at (10.30,2.6) {$T^{\max}$};
 \draw[blue!65!black,very thick] (1.45,1.26)
   .. controls (3.2,1.6) and (3.6,3.32) .. (5.1,3.36)
   .. controls (6.55,3.4) and (7.1,1.42) .. (9.7,1.26);
 \draw[gray!55,densely dotted] (1.45,.45)--(1.45,3.65);
 \draw[gray!55,densely dotted] (9.7,.45)--(9.7,3.65);
 \fill[blue!75!black] (1.45,1.26) circle (.07);
 \fill[blue!75!black] (9.7,1.26) circle (.07);
 \fill[red!70!black] (5.55,3.30) circle (.075);
 \node[font=\scriptsize,anchor=north] at (1.45,.30) {$e_\ell$};
 \node[font=\scriptsize,anchor=north] at (9.7,.30) {$e_{\ell+1}$};
 \draw[-Latex,red!65!black,thick] (6.1,3.86)--(5.63,3.41);
 \node[font=\scriptsize,anchor=west,red!65!black] at (6.13,3.84) {片段内部超温};
 \node[font=\scriptsize,blue!65!black] at (2.9,1.08) {端点安全};
 \node[font=\scriptsize,blue!65!black] at (8.30,1.07) {端点安全};
 \node[font=\scriptsize,align=center] at (5.5,-.58) {端点条件仅是连续温度安全的必要条件};
\end{tikzpicture}}
\caption{事件端点温度合格但片段内部可能超温的概念示意（非实测轨迹）}
\label{fig:ch3-interior-peak}
\end{figure}
